% The local response package constructed at the selected scale.
\section{Construction of the finite local response package}
\label{sec:ex-lattice-trace-eve-producer}

We give the quantitative construction behind the local response package used
by the later trace estimates. The output contains its probe error, chart
error, price, and cloud allowance; these are built together because the final
cloud billing uses the same price budget as the probe estimate.

Fix the scale data of the trace construction and write
\[
 \chi(x)=\sqrt{H(x)},\qquad e=\mathrm{title}(P_1),\quad
 a=\mathrm{title}(P_0),\quad
 h=\mathrm{credibility}(P_1),\quad \nu=\mathrm{DoorsRate}(L).
\]
Here $H$ and $\chi$ are the monitor and its square root from
Section~\ref{sec:ex-lattice-trace-local-rate-selection}, whose functions
$\beta(x)$, $b(x)$, $R(x)$ and $\rho(x)$ are also used below; we write
$b(\beta(x))$ for the radius function evaluated at $\beta(x)$,
\[
 \operatorname{temp}(x)=x^{-1/2}\exp\bigl((\log\log x)^2\bigr),
 \qquad
 \operatorname{fiery}(x)=\operatorname{temp}(\beta(x))
                        =\beta(x)^{-1/2}H(x).
\]
Further, $P_0$ and $P_1$
are the two stage profiles of the transition, and $L$ is the block side;
$e,a,h,\nu$ are positive, $e,a\ge1$, and $\nu=L^{-2/5}$. Let $\mathscr F$ be the grade-$d$ coefficient frame and $F$ its
finite family of coefficient-recovery probes; $F$ is distinct from the
finite basis of the frame, whose indices
are denoted by $\alpha$ below. The prices $\mathrm{testPrice}$,
$\mathrm{readPrice}$, $\mathrm{errorPrice}$ and $\mathrm{orbitCost}$ below
are constants of the frame $\mathscr F$ (the second and third also of the
selected holonomy map), not of the probe family. For $p\in F$, let $w_p=\pi(p)\ge1$ be its
probe price, and write $\Lambda=M_d$ for the polarization constant
in the probe estimate, both as in
Section~\ref{sec:ex-lattice-trace-price-recurrence}. Put
\[
 Z=\mathrm{testPrice}(\mathscr F),\qquad T=\mathrm{readPrice}(\mathscr F),\qquad
 c=\mathrm{orbitCost}(\mathscr F)\,\mathrm{errorPrice}(\mathscr F)\,\mathrm{corn}(d)\,a.
\]
The source normalizations give $Z,T\ge1$ and $c>0$. For a positive target
price $\zeta$, define
\[
 \kappa=\min\left\{\frac{\nu}{2cZ},\zeta\right\}>0,
 \qquad
 g_p(s)=w_p\Lambda L^d\,2(es+s)+s.
 \tag{1}\label{eq:ex-lattice-eve-gp}
\]
The source's finite probe family is essential here: each $g_p$ is continuous,
$g_p(0)=0$, and therefore there is one $b>0$ such that $g_p(s)\le\kappa$
for all $p\in F$ and $0\le s\le b$. Set $Q=b/2$.

\begin{theorem}[Finite-probe and cloud billing]
\label{thm:ex-lattice-eve-producer}
Let $d\in\mathbb N$, and assume $e,a,h,\nu,c,Z,T,\zeta>0$, $e,a,Z,T\ge1$,
and a finite probe family $F$ with weights $w_p\ge0$ and normalization
$\Lambda\ge0$. Suppose the selected response satisfies
$u\in\mathcal R_{-1/2}$ and $\chi\in\mathcal R_0$, with $\chi(x)\ge1$
eventually.
For each probe, the charged scalar is the continuous function
\eqref{eq:ex-lattice-eve-gp}. Set
\[
 \kappa=\min\left\{\frac{\nu}{2cZ},\zeta\right\},\qquad
 \delta_x=Q/\chi(x),\qquad C_x=\kappa/\chi(x).
\]
Assume the following selected-scale receipts hold on one common eventual
tail in $x$; the tail is fixed before the probe, the frame index and all
the other data listed in each receipt are chosen, and each receipt holds
uniformly in those data. For the threshold $\theta=\operatorname{temp}(x)$
and the fine threshold $\Theta=\operatorname{fiery}(x)$,
\[
 \tfrac12\le \frac{\Theta}{\theta}\le2.
 \tag{2}\label{eq:ex-lattice-eve-blend}
\]
The probe receipt is as follows. Let $p\in F$ have order $r_p$ and probe
functional $R_p$ on coefficients, and let $\alpha$ be a basis index of the
coefficient frame, with basis coefficient $\mathsf e_\alpha$. Consider any
step whose block count is $L$ and whose first-order row is literally the
row of the selected kernel with block side $L$, with coarse side larger
than $2\rho(x)$, and any two thresholds $w,y\in(0,1]$ within a factor two
of each other. Let $U_\alpha$ be the coarse-lattice function obtained from the
local presentation of $\mathsf e_\alpha$ at the origin with threshold $w$
by cropping the coarse configuration to radius $\rho(x)$ about the parent
of the origin and routing it through the first-order row of the step; let
$V_p$ be the logarithmic-link test curve of $p$ with amplitude $y$; and
let $\mathcal L$ be the leading coefficient map determined by the selected
kernel and the two thresholds. The
\emph{probe-to-run difference} is
\[
 L^d\left.\frac{d^{r_p}}{ds^{r_p}}U_\alpha(V_p(s))\right|_{s=0}
   -R_p\bigl(\mathcal L\,\mathsf e_\alpha\bigr),
\]
and the receipt is that its absolute value is at most $Q/\chi(x)$ for
every $p\in F$, every $\alpha$, and all such steps and thresholds. The cap
receipt is: for every support $X$ and every pair of budgets $(P,B)$ such
that $P$ has threshold $\theta$ and the power of $P_0$, $B$ has threshold
$\Theta$ and radius $b(\beta(x))$, and
$\operatorname{coverNumber}(B.B,X)\le1$, and for every
$\beta\ge\beta(x)/2$, the scale cap $O_x$ (the response cost of
\eqref{eq:ex-lattice-explicit-local-cost} without its last, paint, line)
and the local moment $M_x$ (the
head-count moment of order $4(n+1)$ there) obey
\[
 0\le O_x\le1,\qquad
 \mathrm{clocks}(e)\,O_x\,M_x\le \frac{Q}{\chi(x)}.
 \tag{3}\label{eq:ex-lattice-eve-cap}
\]
For every cloud $Z_0$ of mass $m_0\le h$, put
$\mathcal M_x(Z_0)=\mathrm{marvel}_x(Z_0)\ge0$. Its cloud receipt is
\[
 L^d\,2\mathcal M_x(Z_0)\le S m_0,
 \qquad S=\frac{\nu}{2cTh}.
 \tag{4}\label{eq:ex-lattice-eve-cloud}
\]
Then for all sufficiently large $x$, with the same uniformity in the data
of each receipt:

\begin{enumerate}
\item the two selected thresholds satisfy the factor-two comparison;
\item the response amplitude $s_x=|u(x)|$ obeys $s_x\le\delta_x$ and
      $s_x\le\zeta/\chi(x)$;
\item the scale cap is at most one, and its clock-weighted moment is at most
      $\delta_x$;
\item every selected probe has approximation error at most $\delta_x$;
\item for each probe $p$, the explicit charge
      $w_p\Lambda L^d\,2(e s_x+\delta_x)+\delta_x$ is at most $C_x$;
\item for every cloud $Z_0$ of mass at most $h$, the combined charge
\[
 c\left[Z C_x+T L^d\,2\mathcal M_x(Z_0)\right]\le\nu.
\]
\end{enumerate}
In particular $C_x\le\zeta/\chi(x)$.
\end{theorem}

\begin{proof}
The product rate gives $\chi u\in\mathcal R_{-1/2}$. Choose
$q\in(-1/2,0)$; the strict-power definition implies
$\chi(x)|u(x)|\le\beta(x)^q\to0$. Thus, eventually,
\begin{equation}
 |u(x)|\le Q/\chi(x),\qquad |u(x)|\le\zeta/\chi(x),\qquad \chi(x)\ge1.
 \label{eq:ex-lattice-eve-response-bound}
\end{equation}
The first two give the asserted response amplitude bound and the desired
price bound for the amplitude. The first selected geometry estimate gives
the factor-two threshold comparison.

The selected scale-cap estimate is called at price $Q/e$. Its price bound
is linear in the allowed scalar multiplier $t$; evaluating it at $t=e$
gives exactly the second inequality in (3), while its companion bound gives
$O_x\le1$. The probe approximation estimate is called at tolerance $Q$,
which gives the bound stated above directly.

For a probe $p$, the amplitude bound
\eqref{eq:ex-lattice-eve-response-bound} gives
\begin{equation}\label{eq:ex-lattice-eve-amplitude}
 2(e|u(x)|+\delta_x)
       \le \frac{2(eQ+Q)}{\chi(x)}.
\end{equation}
Consequently its charged expression is bounded by
\[
 \frac{w_p\Lambda L^d\,2(eQ+Q)+Q}{\chi(x)}
   =\frac{g_p(Q)}{\chi(x)}
   \le \frac{\kappa}{\chi(x)}=C_x.
\]
This is the simultaneous finite-probe choice made above; it is not a
pointwise choice depending on the probe.

Finally, let a cloud have mass $m_0\le h$. The cloud estimate (4) gives
\[
 T L^d(2\mathcal M_x(Z_0))\le T S h=\frac{\nu}{2c}.
\]
The probe part costs at most $ZC_x\le Z\kappa/\chi(x)\le Z\kappa$.
By the definition of $\kappa$, $cZ\kappa\le\nu/2$. Adding the cloud and
probe charges therefore gives at most $\nu$. This proves all six clauses.
\end{proof}

\subsection{Matching the finite-volume response estimate}
\label{subsec:ex-lattice-eve-matching}

We now identify the local insertion receipt in the actual response
construction. Fix a leaf insertion of weight $\tau$, and hold its masked
ingredient fixed while varying the local field. The two functions compared
in the source are the actual-sector and Gaussian-reference integrals
\[
 \Phi_6(y)=\int R_6(y,z)\chi_6(y,z)h(z)\,d\gamma_{C_0}(z),\qquad
 \Phi_9(y)=\int R_9(y,z)\chi_9(y,z)h(z)\,d\gamma_{C_0}(z),
 \tag{5}\label{eq:ex-lattice-eve-two-integrals}
\]
where $h$ is the fixed measurable masked insertion and $|h|\le\tau$.
The density comparison in Section~\ref{sec:principal-gaussian-comparison}
proves, under its displayed finite-dimensional covariance, mean, cutoff,
and chart hypotheses,
\[
 \|\Phi_6-\Phi_9\|_{J^n_0}
 \le \tau P_n(N,Q,S)\bigl(\delta+2N e^{-v+M/2}\bigr).
 \tag{6}\label{eq:ex-lattice-eve-linear-receipt}
\]
The primitive mean and covariance estimates supplying $\delta$ are proved
in Section~\ref{sec:principal-comparison-jets}. The fixed insertion is not
differentiated. The local chart directions are the same ordered variations
used in the response test, with $j\le n$ and their prescribed directional
and canopy bounds.

In the exact construction, the fixed leaf is held together with its regular
masked ingredient $L'$. For each fixed grade and block, the bounds
$N(x),Q(x),S(x),\delta(x)$ and the cutoff-tail bounds are common to every
admissible period, forest and tag choice, every leaf whose cover is at most
$B$, and every normalized direction list allowed in the comparison. They
are independent of the leaf weight $\tau$. Thus the one finite-probe choice
and selected response $u(x)$ work uniformly over all these inputs. The
local read set, forest, tags, insertion mask, and chart are fixed during each
individual differentiation. The two local observables are
\[
 \mathcal O_{\mathrm{act}}(y)=\Phi_6(y),\qquad
 \mathcal O_{\mathrm{ref}}(y)=\Phi_9(y)
 \tag{7}\label{eq:ex-lattice-eve-source-identification}
\]
and these identities follow by writing each side with its own normalized
Gaussian density and using the protected-support agreement on the
insertion's read set. Thus the
response difference is exactly $\Phi_6-\Phi_9$, not a comparison
of different insertions or different Gaussian spaces. The local finite
price cutoff makes the integrands vanish outside a common compact coordinate
ball; the cutoffs are one on the protected smaller ball. Differentiation
under the integral is justified with the bounded measurable $h$ fixed, by
the uniform derivative bounds on that compact ball and dominated
convergence.

The selected rate budgets choose the scalar envelope
\[
 B_*(x)=P_n(N(x),Q(x),S(x))
             \bigl(\delta(x)+2N(x)e^{-v(x)+M(x)/2}\bigr).
\]
The covariance discrepancy has strict-power rate $-1/2$, the mean
discrepancy has rate $-1$ before conversion to Gaussian coordinates, and
the good-coordinate tail has every negative rate. The fixed polynomial
factor has rate zero. Therefore $B_*\in\mathcal R_{-1/2}$.
The response construction packages the concrete derivative, cutoff, and
tail budgets into one scalar envelope. The component bounds are assembled
from the finite product and sum rules for the displayed rates. With this
envelope, for every admissible leaf and every ordered list of at most $n$
normalized chart directions, the derivative of the actual-minus-reference
observable at the chart origin is bounded by $\tau B_*(x)$, uniformly over
the choices above. This is the concrete uniform derivative inequality used
by the response construction.

If the branch patch required by the uniform derivative inequality does not
exist, the response is zero and that inequality is vacuous. If the patch
exists, the source's finite-volume estimate supplies a function $f$ of rate
$-1/2$ and the derivative bound with price $2|f|$. Choosing the response to
be $u=2|f|$ preserves that rate and bound. Thus the remaining assembly in
this section is exactly the finite-test choice and scalar billing proved
above; the local insertion estimate is supplied by the two cited comparison
sections under their explicit hypotheses.

For the record, the source constructs the package for one transition with
block side $L$ and grade $d$ under the following standing conditions, which
the scalar theorem above does not see: the guard of the record satisfies
$D.\varepsilon\le1/4$; the two integer parameters of the selected kernel
are at least $4$, and the first of them is at least $d$; the differentiation
order satisfies $n\ge2$; and the real parameter $\mu$ of the response
estimate is positive. The response receipt with amplitude $|u(x)|$ is stated there for
every leaf whose label is the logarithmic key of a regular activity and
whose ingredient is exactly the mask, at the threshold of the
representation, of a function that is regular on the support of that
activity; for every budget of radius $b(x)$ and of the power of $P_0$ that
is compatible with the representation; for every read support that is
healthy at the base configuration and has cover number at most that
radius; and for every list of at most $n$ directions satisfying the
budget's directional bound and the canopy bound
$2\Theta(1+R(x))^{p}$, where $p$ is the power of $P_0$. These identities
of the leaf, mask, budget and support are what make
\eqref{eq:ex-lattice-eve-source-identification} a comparison of one
insertion on one Gaussian space. The cap receipt of
Theorem~\ref{thm:ex-lattice-eve-producer} follows from the scalar budget
estimate \eqref{eq:ex-lattice-budget-output} of
Section~\ref{sec:ex-lattice-trace-local-rate-selection}, called
at price $Q/e$ and with the nonnegative paint term omitted. The probe
receipt and the cloud receipt (4) are separate inputs, not proved in this
section.
